\summarychapter{41}{Nothing to Lose but Chains}{The Social Crisis}

\begin{inbrief}
The Industrial Revolution multiplied wealth and crowded workers into factories and slums. It set
before Christianity a challenge to match the conflict with modern thought. The chapter describes
factory life and the \emph{laissez faire} creed, the rise of socialism and Karl Marx, and the
Christian answers: Catholic social teaching, the English Nonconformists, the Salvation Army, the
Christian Socialists and the Social Gospel in America.
\question{How did Christians respond to the social crisis created by the industrial city?}
\end{inbrief}

In \emph{Hard Times} Charles Dickens drew a typical English factory town: brick that would have been
red if smoke and ashes had allowed it, tall chimneys trailing ``interminable serpents of smoke,'' a
river running purple with ill-smelling dye, and vast buildings rattling and trembling all day.
Historians call the rapid changes of the nineteenth century the industrial revolution because most of
them go back to new methods of manufacture. Factories meant cities, cities meant people, and people
meant problems. Many critics still insist that Christianity never faced the crisis but fled from it.
Large territories once shaped by Christian values came under Marxist governments, and even in the
West few people see any connection between their religion and their work.

\ssection{The New World of Factories}
Many believers find the industrial city a riddle and fear its sins and crimes. Yet Christianity has
always been both a conservative and a radical force, and it was both here. France had a political
revolution and Germany an intellectual one, but the industrial revolution was born in England, which
as the leading commercial and financial power had the markets and the money. Of the many
breakthroughs none mattered more than steam. James Watt (1736--1819), the Scottish genius at the
University of Glasgow, found how to control the supply of steam and drive a piston in a closed
cylinder, and steam engines soon ran textile mills and later locomotives and ships. For the first
time manufacturers had power that did not depend on weather or season. The factory became the symbol
of the new order, and people left cottages, looms and fields to tend its machines. Rural life, paced
by night and day and by sowing and harvest, gave way to precision and regimen. God's sun was hidden
by smoke, and in its place sounded the factory whistle, the sign of man's time, not God's.

The new wealth brought a host of evils for the workers. Early factories lacked the simplest sanitary
and safety provisions, and mangling accidents were common. Under English common law an accident was
put down to the worker's own negligence, so the employer was not liable; with no compensation or
health insurance, the injured were turned out destitute and replaced from the thousands seeking work.
Women and children worked beside men. Poverty kept women at work until a day or two before
childbirth and brought them back soon after, and many mills employed children of four or five. The
day was twelve to fifteen hours of hard labor for minimum wages. A family lived in one room of a
dirty tenement, lucky not to share it, among streets strewn with garbage and, for want of sewers,
reeking of excrement.

The old conflict between aristocrats and middle class had ended in the triumph of the bourgeoisie,
and with it came the philosophy of \emph{laissez faire}. It held that the social ills of
industrialism could not be corrected, and that if each individual is left alone to pursue his own
interest, the result will be the happiness of the greatest number. Its textbook, Adam Smith's
\emph{Wealth of Nations}, argued that society gains most from competition, which rewards the
efficient. Government should protect life and property and keep its hands off business, leaving the
laws of supply and demand to work. The urban workers had no share in the benefits. Having no
property they had no vote, and their struggle in the nineteenth century was for the ballot or some
other means of power.

The first attack on \emph{laissez faire} came from socialism. At this stage it was not a militant
workers' movement but a theory that condemned the concentration of wealth, called for public or
worker ownership, and insisted that harmony and cooperation, not ruthless competition, should govern
economic life. These early socialists are called utopians because their model communities rested on
the naive belief that people naturally love one another and that only capitalist competition sets
them against each other. The problems of industrial society thus came down to one question, property.
The church had always affirmed a natural right to own things, which seemed embodied in the tenth
commandment, and had treated those who denied it as heretics. But never had there been so much
property or so few allowed to own it. Millions, not thousands, were wholly dispossessed. Rome had
appeased its people with bread and circuses; now the poor had neither. Hunger filled the minds of the
lowest classes and insecurity crept among the middle classes.

The churches were poorly placed to speak. The separation of church and state had left them without a
political platform. The American Constitution forbade a national church, the English reforms of the
1830s stripped the Church of England of many privileges, and in Catholic lands such as France and
Italy anticlerical feeling limited the hierarchy. The churches were shut out of the arena where the
crucial decisions were made. People inside and outside them came to think of Christianity in ever
narrower terms, leaving whatever was not plainly ``spiritual'' beyond criticism, and a growing body
of industrial workers judged the Christian message irrelevant or powerless in the age of the
machine.

\ssection{Marx, Prophet of the New Age}
Almost unnoticed, a prophet arose in 1848 with a new form of socialism. Karl Marx charged
Christianity with betraying its own vision by siding with the wealthy oppressors. Like religion in
general it was a tool of the rich. It told the poor to bear injustice now for a reward in heaven and
distracted them with lofty talk of noble ideas, whereas present conditions are in fact the result of
conflict between social classes. The time had come for the poor to rise in history-changing
violence. Marx (1818--1883) was born at Trier in the Rhineland to German-Jewish parents who had
become Christians, and took his doctorate after studying the philosophy of Hegel. Finding no
university post, he lived precariously as a journalist. In Paris he took up socialist ideas and began
his lifelong friendship with Friedrich Engels (1820--1895), the son of a rich German factory owner.
Expelled by the French in 1845, he went with Engels to Brussels.

In January 1848 the two published the \emph{Communist Manifesto}, which contains nearly all the
elements of what they called ``scientific'' socialism. It opens, ``A spectre is haunting
Europe---the spectre of Communism,'' calls for implacable struggle against the bourgeoisie, and
proclaims the inevitable revolution and the triumph of the masses. It ends with the warning that
communist aims can be reached only by the forcible overthrow of all existing social conditions:
``The proletarians have nothing to lose but their chains. They have a world to win. Working men of
all countries, unite!'' After a short time in Germany Marx spent the rest of his life in London,
mostly on the edge of survival and supported by friends, above all Engels. Almost daily he went to
the British Museum to gather material for his books, especially \emph{Das Kapital} (1867--1894), one
of the most influential books of modern times. For a season in the twentieth century nearly half the
world lived under governments rooted in Marxism.

As in the slavery controversy, Christians stood on every side. Many church members were capitalists,
owners of factories with influence in society and eager to defend \emph{laissez faire}. After
mid-century, however, more and more Christians, Catholic and Protestant, worked for better conditions
for laborers. Four courses were open to them: (1) to challenge \emph{laissez faire} in the name of
Christian principles; (2) to found Christian institutions for the relief of the poor and powerless;
(3) to support the formation of labor unions; and (4) to appeal to the state for laws improving
working conditions.

Among Catholics the German bishop Wilhelm Ketteler (1811--1877) addressed the factory worker's plight
in sermons and books as early as 1848. He saw danger both in the unlimited competition of capitalism
and in the exaggerated state control of the socialists. Against the capitalists he defended the
state's right to intervene, and against the socialists the right of private property. Above all he
insisted on the workers' right to form their own associations, and he called for profit sharing,
reasonable hours, enough days of rest, factory inspection and the regulation of women's and
children's labor. In England Cardinal Henry Edward Manning (d.~1892), a convert from evangelical
Anglicanism, appeared in December 1872 at a meeting in support of farm workers, the first time a
Catholic prelate there had so openly taken labor's side. He urged Gladstone to prohibit the labor of
children under a certain age and to regulate housing, and in his 1874 lecture ``The Rights and
Dignity of Labour'' he defended the right to organize, called for laws on working hours and exposed
the abuses of child labor. In the United States the Catholic layman Terence Powderly led the first
really effective union, the Knights of Labor, open to all but with a large Catholic membership from
the start. Roman officials, the pope included, feared unions, remembering their conflicts with the
Italian Masons and the forces of Italian unity. The Knights probably escaped condemnation only
because the archbishop of Baltimore, later Cardinal Gibbons, defended them.

The papacy spoke only in 1891. The encyclical \emph{Rerum Novarum} of the aging Leo XIII
(1878--1903) set out what became the basic Catholic position on capital and labor. After describing
the evils that came from the industrial revolution, it rejects socialism and upholds the Christian
family as the most essential unit of society. But it also stresses the capitalist's duty to deal
justly with workers, and defends the right of Catholic workers to organize and bargain and their
right to a living wage and reasonable leisure. Monopoly capitalism and industrial slavery thus came
under papal criticism along with socialism.

\ssection{Laboring Men in England}
In largely Protestant England support for the worker came chiefly from the Nonconformists. Unions
were illegal until mid-century, but the labor movement found strong backing among them, especially
the Primitive Methodists. Many labor leaders had served their apprenticeship as Methodist lay
preachers and carried into the workingmen's associations the zeal, the organizing skill and the
speaking experience of the chapel. Nonconformists also led the temperance crusade and founded
orphanages: George M\"uller of the Plymouth Brethren, Charles Haddon Spurgeon, London's leading
Baptist preacher, and the Methodist T.~B.~Stephenson.

\ssection{Onward, Christian Soldiers!}
The outstanding ministry to the dispossessed was that of the pietistic Evangelical William Booth
(1829--1912). He began with the Methodist New Connection but soon withdrew to work among London's
poor. His street preaching in the East End in 1864 met extraordinary success, and within eleven
years he had thirty-two stations for evangelism and social service. His workers, organized like a
military unit, were soon called the Salvation Army, and the evangelist became General Booth. By 1888
he had a thousand British corps and patrols in many other nations. His book \emph{In Darkest England
and the Way Out} (1890) compared England's social darkness with the African darkness described by
Livingstone. In one year in London, he reported, 2,157 people had been found dead and 2,297 had
killed themselves; 30,000 lived by prostitution, 160,000 had been convicted of drunkenness, and more
than 900,000 were classed as paupers. He went on to describe the Army's vast rescue work.

The Church of England, wedded to the past, adjusted with difficulty. A new parish in a mushrooming
industrial town required an act of Parliament, slow and costly, so the urban masses usually grew up
outside its care. The striking exception was the Christian Socialist group of the theologian
F.~D.~Maurice (1805--1872), the novelist Charles Kingsley (1819--1875) and the lawyer John Malcolm
Ludlow (1821--1911), who believed the gospel had something to say to working people. In its brief
active period, 1848 to 1854, the movement attacked the whole \emph{laissez faire} outlook.
``Competition is put forth as the law of the universe,'' Maurice wrote to Kingsley. ``That is a
lie.'' The true law is that man is made to live in community, and people realize their nature when
they cooperate as children of God and brothers in Christ. The leaders took the name socialism to be
provocative. By it they meant no more than the science of making men partners, and they regarded it
as a development of Christianity. The practical results were small. Their cooperative workshops were
badly organized and too optimistic about the workers, and their main contribution was probably the
ideas that crossed to the United States. English workers meanwhile gained political power and better
conditions step by step, as act after act whittled away \emph{laissez faire}: no work for children
under ten, a ten-hour day for women and children, safety inspection of factories. England had done
much for its laborers before many Americans recognized that a social crisis existed.

\ssection{Social Gospel in America}
The chief American movement for social justice was the Social Gospel, the work of liberal Protestant
pastors and professors who formed no organization of their own but worked through the denominations
and political agencies. Its crux was the belief that God's saving work includes corporate structures
as well as personal lives. If social good and evil are collective, not merely the sum of good and
evil individuals, then Christians are bound to work for the reconstruction of the social order. Its
prophets drew on several sources. Revivalism, heavily individualistic under Moody, had once carried
a vision of a moral, Christian America, and before the Civil War revival and reform went together;
Finney urged his converts on from regeneration to the antislavery crusade. They also read the
European Christian Socialists, Maurice and Kingsley among them. And the New Theology, the
``progressive orthodoxy'' of Andover and the New England pulpits, supplied their general pattern of
belief.

The father of the movement was Washington Gladden (1836--1918), a transplanted New Englander who
published the first of his many books on the subject in 1876 and spent his most influential years at
the First Congregational Church of Columbus, Ohio. With employers and employees both in his
congregation, he watched with alarm ``the widening of the breach between these classes'' in times of
strife, and in evening addresses on the labor problem he argued that the teachings of Jesus contain
the principles for the right ordering of society. His spirit is in his hymn ``O Master, let me walk
with Thee / In lonely paths of service free.'' He was no socialist. He kept private property and
private enterprise, yet thought many industries could be run cooperatively and that railroads, mines
and city utilities should be operated by government.

The prophet of most lasting influence was Walter Rauschenbusch (1861--1918). As a young German
Baptist pastor in Hell's Kitchen, a tenement district of New York, he wrestled with the Christian
response to urban problems, and as professor of church history at Rochester Theological Seminary he
wrote \emph{Christianity and the Social Crisis} (1907), \emph{Christianizing the Social Order} (1912)
and \emph{A Theology for the Social Gospel} (1917). He avoided the comfortable doctrine of progress
and anchored social responsibility in the kingdom of God. The social gospel, he wrote, is ``the old
message of salvation, but enlarged and intensified.'' The individualistic gospel has shown the
sinfulness of every heart and God's power to save every soul that comes to him, but not the
sinfulness of the social order and its share in the sins of all within it; the social gospel seeks to
bring people to repentance for their collective sins and to create a more sensitive and more modern
conscience. The prime example of social sin, its preachers said, was the capitalist system, and
salvation was impossible while it stayed unchanged. They differed about how much change was needed
but agreed that the kingdom of God could not come without it.

Many seminaries reshaped their curricula to address social concerns, but the chief sign of the
movement's reach was the formation of the Federal Council of Churches in 1908. Almost its first act
was to adopt a ``Social Creed of the Churches,'' calling for benefits later American workers took for
granted: occupational safety, security in old age, a minimum wage and the right of arbitration.
Throughout its history the church has tried to better human life on earth while preparing people for
the world to come, and it has always been possible to dwell on the next life so much that Christians
seemed insensitive to present pain. The movements for social concern ran the opposite risk of
reducing the gospel to social activism. Their lasting reminder, Shelley concludes, is that Christians
cannot show concern for people's eternal destiny unless they also show concern for their earthly
needs.

\begin{keyfacts}
\keypeople{James Watt; Adam Smith; Karl Marx; Friedrich Engels; Wilhelm Ketteler; Cardinal Manning;
Terence Powderly; Cardinal Gibbons; Leo XIII; William Booth; F.~D.~Maurice; Charles Kingsley;
Washington Gladden; Walter Rauschenbusch.}
\keydates{%
  \kd{1848}{Marx and Engels publish the \emph{Communist Manifesto}}
  \kd{1848--1854}{Christian Socialist movement in England}
  \kd{1864}{William Booth begins preaching in London's East End}
  \kd{1867}{First volume of Marx's \emph{Das Kapital}}
  \kd{1890}{Booth's \emph{In Darkest England and the Way Out}}
  \kd{1891}{Leo XIII's encyclical \emph{Rerum Novarum}}
  \kd{1907}{Rauschenbusch's \emph{Christianity and the Social Crisis}}
  \kd{1908}{Federal Council of Churches adopts the ``Social Creed of the Churches''}}
\keyterms{Industrial Revolution; \emph{laissez faire}; bourgeoisie; proletariat; utopian socialism;
``scientific'' socialism; communism; labor unions; Salvation Army; Christian Socialism; Social
Gospel; kingdom of God; social sin.}
\end{keyfacts}
